
% this file is called up by thesis.tex
% content in this file will be fed into the main document
%----------------------- introduction file header -----------------------
%%%%%%%%%%%%%%%%%%%%%%%%%%%%%%%%%%%%%%%%%%%%%%%%%%%%%%%%%%%%%%%%%%%%%%%%%
%  Capítulo 1: Introducción- DEFINIR OBJETIVOS DE LA TESIS              %
%%%%%%%%%%%%%%%%%%%%%%%%%%%%%%%%%%%%%%%%%%%%%%%%%%%%%%%%%%%%%%%%%%%%%%%%%

\chapter{Introducción}

%: ----------------------- HELP: latex document organisation
% the commands below help you to subdivide and organise your thesis
%    \chapter{}       = level 1, top level
%    \section{}       = level 2
%    \subsection{}    = level 3
%    \subsubsection{} = level 4
%%%%%%%%%%%%%%%%%%%%%%%%%%%%%%%%%%%%%%%%%%%%%%%%%%%%%%%%%%%%%%%%%%%%%%%%%
%                           Presentación                                %
%%%%%%%%%%%%%%%%%%%%%%%%%%%%%%%%%%%%%%%%%%%%%%%%%%%%%%%%%%%%%%%%%%%%%%%%%

\section{Presentación} % section headings are printed smaller than chapter names


%%%%%%%%%%%%%%%%%%%%%%%%%%%%%%%%%%%%%%%%%%%%%%%%%%%%%%%%%%%%%%%%%%%%%%%%%
%                           Objetivo                                    %
%%%%%%%%%%%%%%%%%%%%%%%%%%%%%%%%%%%%%%%%%%%%%%%%%%%%%%%%%%%%%%%%%%%%%%%%%

\section{Objetivo}

El principal objetivo de este trabajo es desarrollar, analizar y modelar el comportamiento de los clientes de cualquier empresa en México, el enfoque de este trabajo puede ser dirigido tanto a empresas que ofrezcan servicios de retail, suscripciones, telecomunicaciones, financieros/bancarios o inclusive empresas de nicho como las pertenecientes al sector agrícola, todo esto utilizando técnicas de aprendizaje automático, más en especifico usaremos modelos de clasificación para predecir su comportamiento, la idea es identificar patrones de conducta asociados a la retención y al abandono de clientes (Churn), haciendo uso de herramientas de ciencia de datos implementadas en Python. La finalidad es permitir a cualquier tipo de empresas anticipar las posibles perdidas de clientes y diseñar diversas estrategias de fidelización.
%%%%%%%%%%%%%%%%%%%%%%%%%%%%%%%%%%%%%%%%%%%%%%%%%%%%%%%%%%%%%%%%%%%%%%%%%
%                           Motivación y estado del arte                %
%%%%%%%%%%%%%%%%%%%%%%%%%%%%%%%%%%%%%%%%%%%%%%%%%%%%%%%%%%%%%%%%%%%%%%%%%
\section{Motivación}

Popularmente se cree que adquirir nuevos clientes cuesta más que mantener los actuales, tal afirmación plantea la importancia que toma el análisis de churn en la toma de decisiones. Hoy en día, la retención de clientes representa uno de los mayores desafíos para las empresas en todo tipo de sectores, este caso de estudio se ha aplicado a empresas de retail, telecomunicaciones, servicios financieros, entre otros. La innovación científica y tecnológica en las empresas de México y todo el mundo avanza a pasos agigantados y con ella surgen distintas necesidades, el poder incorporar herramientas de análisis de datos para poder comprender mejor el comportamiento de sus clientes es una de ellas, en especial a la hora de comprender las razones por las cuales sus clientes dejan de consumir sus servicios/productos, por lo cual el análisis de abandono de clientes se ha convertido en un requisito  estratégico. Esta investigación surge por la motivación de integrar un enfoque cuantitativo, dirigido al Machine Learning, esto con un enfoque científico para poder ofrecer una solución innovadora a todo tipo sectores, incluso aquellos donde la digitalización aun esta en desarrollo.

%%%%%%%%%%%%%%%%%%%%%%%%%%%%%%%%%%%%%%%%%%%%%%%%%%%%%%%%%%%%%%%%%%%%%%%%%
%                   Planteamiento del problema                          %
%%%%%%%%%%%%%%%%%%%%%%%%%%%%%%%%%%%%%%%%%%%%%%%%%%%%%%%%%%%%%%%%%%%%%%%%%

\section{Planteamiento del problema}

El objetivo será desarrollar un modelo de explicativo y predictivo del Churn utilizando diferentes algoritmos de Machine Learning, comparando tanto modelos clásicos (por ejemplo regresión logística, maquinas de soporte vectorial, entre otros) así como modelos más complejos (XGBoost, redes neuronales, etc.), de tal manera que podamos estimar la capacidad de estos modelos para predecir el Churn y explicar el porqué ciertos clientes tienen más probabilidad de abandonar. La importancia de ofrecer una buena alternativa que explique y prediga dicho fenómeno ayudaría a poder generar hallazgos y accionables.

%%%%%%%%%%%%%%%%%%%%%%%%%%%%%%%%%%%%%%%%%%%%%%%%%%%%%%%%%%%%%%%%%%%%%%%%%
% Preguntas de  investigación                                
%%%%%%%%%%%%%%%%%%%%%%%%%%%%%%%%%%%%%%%%%%%%%%%%%%%%%%%%%%%%%%%%%%%%%%%%%

\section{Preguntas de investigación}


%%%%%%%%%%%%%%%%%%%%%%%%%%%%%%%%%%%%%%%%%%%%%%%%%%%%%%%%%%%%%%%%%%%%%%%%%
% Hipótesis o supuestos del estudio                              
%%%%%%%%%%%%%%%%%%%%%%%%%%%%%%%%%%%%%%%%%%%%%%%%%%%%%%%%%%%%%%%%%%%%%%%%%

\section{Hipótesis o supuestos del estudio}

%%%%%%%%%%%%%%%%%%%%%%%%%%%%%%%%%%%%%%%%%%%%%%%%%%%%%%%%%%%%%%%%%%%%%%%%%
%                           Metodología                                 %
%%%%%%%%%%%%%%%%%%%%%%%%%%%%%%%%%%%%%%%%%%%%%%%%%%%%%%%%%%%%%%%%%%%%%%%%%
\section{Metodología}

La metodología propuesta para el desarrollo de este trabajo esta constituida por cuatro etapas primordiales, \textbf{primero} recopilación y preprocesamiento de los datos, \textbf{segundo} implementación de modelos de aprendizaje automático (Machine Learning y Deep Leaning), \textbf{tercero} aplicación de técnicas de inteligencia artificial explicable (XAI) y finalmente en \textbf{cuarto} lugar la evaluación comparativa de los modelos. Dicha estructura nos va a permitir ofrecer un análisis predictivo y explicativo del Churn, de tal suerte que la compresión e interpretación de los factores que explican el abandono de clientes se puedan aplicar de manera trasversal a diferentes sectores de la iniciativa privada. 


%%%%%%%%%%%%%%%%%%%%%%%%%%%%%%%%%%%%%%%%%%%%%%%%%%%%%%%%%%%%%%%%%%%%%%%%%
%                         Contribuciones                                %
%%%%%%%%%%%%%%%%%%%%%%%%%%%%%%%%%%%%%%%%%%%%%%%%%%%%%%%%%%%%%%%%%%%%%%%%%

\section{Contribuciones}

Este trabajo plantea una serie de aportaciones tanto teóricas como prácticas en el campo de la ciencia de datos enfocados al Churn. Las principales aportaciones se describen a continuación:

\begin{enumerate}
	\item Diseño de un pipeline generalizado y modular para la predicción del churn mediante aprendizaje automático.
	\item Comparación sistemática entre modelos clásicos y avanzados de Machine Learning, integrando criterios de precisión y explicabilidad.
	\item Aplicación práctica orientada a la toma de decisiones empresariales.
	\item Contribución al debate sobre la explicabilidad en sistemas de inteligencia artificial aplicados a negocios.
\end{enumerate}



%%%%%%%%%%%%%%%%%%%%%%%%%%%%%%%%%%%%%%%%%%%%%%%%%%%%%%%%%%%%%%%%%%%%%%%%%
%                           Estructura de la tesis                      %
%%%%%%%%%%%%%%%%%%%%%%%%%%%%%%%%%%%%%%%%%%%%%%%%%%%%%%%%%%%%%%%%%%%%%%%%%

\section{Estructura de la tesis}

Este trabajo está dividido en 5 capítulos, todos y cada uno de ellos responde a una fase del proceso de esta investigación, las cuales van desde la contextualización del problema (Churn) hasta el análisis de los resultados y las conclusiones. A continuación, describiremos brevemente el contenido de cada capitulo:

\begin{enumerate}
	\item Introducción
	\item Revisión de la literatura
	\item Metodología
	\item Resultados
	\item Conclusiones
\end{enumerate}